\subsection{From the full electron-nuclear Schr\"odinger equation to Kohn-Sham DFT}
\label{subsec:dft_background}

At the most fundamental nonrelativistic level, an atomistic system is described by
the stationary many-body Schr\"odinger equation for all electrons and all nuclei,
\begin{equation}
\hat{H}(\bm{r};\bm{R}) \Psi(\bm{r};\bm{R})
=
E \Psi(\bm{r};\bm{R}),
\end{equation}
where $\Psi(\bm{r};\bm{R})$ denotes the many-body wavefunction, $\bm{r}$ the
collection of electronic coordinates, and $\bm{R}$ the collection of ionic
coordinates. The semicolon indicates that the dependence on ionic coordinates
is parametric within the Born-Oppenheimer approximation. We adopt atomic units throughout,
with $\hbar = m_e = e^2 = 1$, so that energies are measured in Hartree and
lengths in Bohr radii.

Under the Born-Oppenheimer approximation \cite{born_zur_1927}, the electronic
Hamiltonian is written as
\begin{equation}
=
\hat{T}_e(\bm{r}) + \hat{V}_{ee}(\bm{r}) + \hat{V}_{ei}(\bm{r};\bm{R}) + E_{ii}(\bm{R}),
\end{equation}
with
\begin{equation}
\hat{T}_e(\bm{r}) = -\frac{1}{2}\sum_{j=1}^{N_e} \nabla_j^2,
\end{equation}
\begin{equation}
\hat{V}_{ee}(\bm{r}) =
\frac{1}{2}\sum_{j=1}^{N_e}\sum_{k \ne j}^{N_e}\frac{1}{|\bm{r}_j-\bm{r}_k|},
\end{equation}
\begin{equation}
\hat{V}_{ei}(\bm{r};\bm{R}) =
-\sum_{j=1}^{N_e}\sum_{\alpha=1}^{N_i}\frac{Z_{\alpha}}{|\bm{r}_j-\bm{R}_{\alpha}|},
\end{equation}
and
\begin{equation}
E_{ii}(\bm{R}) =
\frac{1}{2}\sum_{\alpha=1}^{N_i}\sum_{\beta \ne \alpha}^{N_i}
\frac{Z_{\alpha} Z_{\beta}}{|\bm{R}_{\alpha}-\bm{R}_{\beta}|}.
\end{equation}
The ion-ion term amounts to a constant energy shift for fixed ionic
configuration, while the Hamiltonian and wavefunction depend only
parametrically on $\bm{R}$. In the following, this parametric dependence is
not always written explicitly.

The central remaining difficulty is therefore the interacting many-electron
problem defined by $\hat{T}_e + \hat{V}_{ee} + \hat{V}_{ei}$ at fixed ionic
geometry.

Density functional theory (DFT) replaces the many-electron wavefunction by the
electron density as the central variable \cite{HoKo1964}. In the Kohn-Sham
construction \cite{KoSh1965}, the interacting electron problem is mapped onto a
fictitious system of non-interacting electrons governed by
\begin{equation}
\left[
-\frac{1}{2}\nabla^2 + v_S(\bm{r})
\right]\psi_j(\bm{r}) = \epsilon_j \psi_j(\bm{r}),
\end{equation}
where $j=1,\dots,N_e$ labels the Kohn-Sham orbitals. The corresponding
electronic density is
\begin{equation}
n(\bm{r}) = \sum_{j=1}^{N_e} |\psi_j(\bm{r})|^2.
\end{equation}
The Kohn-Sham potential is
\begin{equation}
v_S(\bm{r}) =
\frac{\delta E_H[n]}{\delta n(\bm{r})}
+ \frac{\delta E_{XC}[n]}{\delta n(\bm{r})}
+ v_{ei}(\bm{r}),
\end{equation}
where $E_H[n]$ is the Hartree energy, $E_{XC}[n]$ the exchange-correlation
energy, and
\begin{equation}
v_{ei}(\bm{r}) = -\sum_{\alpha} \frac{Z_{\alpha}}{|\bm{r}-\bm{R}_{\alpha}|}.
\end{equation}
Although the Kohn-Sham construction is formally exact, practical calculations
rely on approximations to $E_{XC}$.

In this notation, the ground-state energy functional is written as
\begin{equation}
E[n] =
T_S[n] + E_H[n] + E_{XC}[n]
+ \int d\bm{r}\, n(\bm{r}) v_{ei}(\bm{r})
+ E_{ii},
\end{equation}
where $T_S[n]$ is the Kohn-Sham kinetic energy. This is the DFT-level starting
point that Mandala connects to localized-orbital operator learning.

Mandala operates on localized-orbital matrix representations of this problem.
Expanding the Kohn-Sham orbitals in a finite basis $\{\chi_{\mu}\}$,
\begin{equation}
\psi_j(\bm{r}) = \sum_{\mu} C_{\mu j}\chi_{\mu}(\bm{r}),
\end{equation}
leads to the generalized eigenvalue problem
\begin{equation}
\sum_{\nu} H_{\mu \nu} C_{\nu j}
=
\epsilon_j \sum_{\nu} S_{\mu \nu} C_{\nu j},
\end{equation}
with
\begin{equation}
H_{\mu \nu} = \langle \chi_{\mu} | \hat{H}_{\mathrm{KS}} | \chi_{\nu} \rangle,
\qquad
S_{\mu \nu} = \langle \chi_{\mu} | \chi_{\nu} \rangle.
\end{equation}
The one-particle density matrix is
\begin{equation}
D_{\mu \nu} = \sum_j f_j C_{\mu j} C_{\nu j}^{*},
\end{equation}
where $f_j$ denotes the orbital occupation. In a nonorthogonal basis, the
electron count is
\begin{equation}
N_e = \Tr(DS),
\end{equation}
while the energy observable used throughout Mandala is obtained from
\begin{equation}
E_{\mathrm{op}} = \Tr(DH).
\end{equation}
For backend-specific total energies, additional double-counting and ion-ion
terms may be required, but the trace expressions above are the core operator
identities used throughout Mandala.

Localized basis sets also induce block sparsity. If orbital $\mu$ belongs to
atom $i$ and orbital $\nu$ belongs to atom $j$, then matrix entries can be
organized into atom-pair blocks $X_{ij}^{\bm{L}}$, where
$X \in \{H,S,D\}$ and $\bm{L} \in \mathbb{Z}^3$ indexes periodic lattice
images. This sparse block structure is the operator domain learned by Mandala.

\subsubsection{Hamiltonian gauge invariance and the overlap-shift correction}
\label{subsubsec:gauge_invariance}

In a nonorthogonal basis, the generalized eigenproblem is invariant under a
rigid shift of the Hamiltonian by a multiple of the overlap matrix. Let
\begin{equation}
H C = S C \varepsilon,
\end{equation}
and define a shifted Hamiltonian
\begin{equation}
\widetilde{H} = H - \alpha S.
\end{equation}
Then
\begin{equation}
\widetilde{H} C
=
(H-\alpha S) C
=
S C \varepsilon - \alpha S C
=
S C (\varepsilon - \alpha I).
\end{equation}
Therefore the generalized eigenvectors are unchanged, while all eigenvalues are
shifted by the same constant $\alpha$. This proves that the choice of zero of
energy can be represented, in matrix form, as the transformation
$H \mapsto H - \alpha S$. For a fixed electron number and a geometry-independent
shift $\alpha$, one also finds
\begin{equation}
\Tr(D\widetilde{H}) = \Tr(DH) - \alpha \Tr(DS) = \Tr(DH) - \alpha N_e,
\end{equation}
so only the absolute energy reference changes.

Mandala exploits this gauge freedom through the overlap-shift correction
$\mu_H$, which is chosen to minimize the mean-squared discrepancy between a
predicted Hamiltonian and a reference Hamiltonian after subtraction of a
multiple of the overlap matrix:
\begin{equation}
\mu_H
=
\arg\min_{\mu}
\left\|
H^{\mathrm{pred}} - H^{\mathrm{ref}} - \mu S
\right\|_F^2.
\end{equation}
Taking the derivative with respect to $\mu$ and setting it to zero yields
\begin{equation}
\mu_H
=
\frac{
\left\langle
H^{\mathrm{pred}} - H^{\mathrm{ref}}, S
\right\rangle_F
}{
\left\langle S, S \right\rangle_F
},
\end{equation}
where $\langle A,B\rangle_F = \sum_{\mu \nu} A_{\mu \nu} B_{\mu \nu}$ is the
Frobenius inner product, evaluated in Mandala blockwise over the aligned sparse
representation. The corrected Hamiltonian
$H^{\mathrm{pred}} - \mu_H S$ therefore removes a pure gauge offset without
altering the physically relevant eigenvectors.
